% GENERATED FILE. Edit canonical lessons, then run npm run generate:books.
% canonical-source-hash: fnv1a64:f7a048c435c26718
% canonical-lessons: MR-C379-maryada, MR-C379-kiman, MR-C379-kamal, MR-C379-sithil-karne, MR-C379-jhepne

\chapter{Limits}
\label{ch:mr-limits}
\hlchaptermodality{379}

\begin{chapteropening}
\textbf{By the end of this chapter:} \emph{I can talk about limits, minimums and maximums.}

Complete the last lesson of chapter 379: I can talk about limits, minimums and maximums.
\end{chapteropening}

\section[\texorpdfstring{\mr{मर्यादा} (maryādā)}{मर्यादा (maryādā)}]{\mr{मर्यादा} (maryādā) — a limit}

\label{lesson:MR-C379-maryada}



\begin{quote}
Before the new one: say the Marathi for a criterion, then the Marathi for to fulfil.
\end{quote}

\subsection*{You'll want to know: \mr{मर्यादा}}
\textbf{\mr{मर्यादा}} — \emph{maryādā} — \textquotedblleft{}a limit\textquotedblright{}.

\begin{quote}
\mr{प्रत्येक} \mr{गोष्टीला} \mr{एक} \mr{मर्यादा} \mr{असते}.
\end{quote}

\begin{quote}
\emph{pratyek goṣṭīlā ek maryādā aste.}
\end{quote}

\begin{quote}
Everything has a limit.
\end{quote}

Feminine.

\mr{मर्यादा} \mr{घालणे}, to set a limit; \mr{मर्यादा} \mr{ओलांडणे}, to overstep it.

\begin{cousinweb}
\textbf{\mr{मर्यादा}} (\emph{maryādā}) is a tatsama, taken from Sanskrit \emph{maryādā}, \textquotedblleft{}region, boundary\textquotedblright{} (CDIAL 9895).
\end{cousinweb}

\begin{grammarlens}[title={What you've built}]
One more word for saying what happens if.
\end{grammarlens}

\subsection*{Your turn}
\begin{itemize}
\raggedright
  \item \emph{Say it:} \emph{maryādā}
  \item \emph{Say it:} \emph{maryādā}, once more
  \item \emph{Say it:} \emph{pratyek goṣṭīlā ek maryādā aste.} — Everything has a limit.
  \item \emph{Recall:} say the Marathi for a criterion, then the Marathi for to fulfil, then say \emph{maryādā} again
\end{itemize}

\begin{tcolorbox}[breakable,colback=teal!4,colframe=teal!35!black,title={Before you move on}]
What does \mr{मर्यादा} mean? (\textquotedblleft{}A limit\textquotedblright{}.) Say it once more.
\end{tcolorbox}
\section[\texorpdfstring{\mr{किमान} (kimān)}{किमान (kimān)}]{\mr{किमान} (kimān) — minimum; at least}

\label{lesson:MR-C379-kiman}



\begin{quote}
Before the new one: say the Marathi for to fulfil, then the Marathi for a limit.
\end{quote}

\subsection*{You'll want to know: \mr{किमान}}
\textbf{\mr{किमान}} — \emph{kimān} — \textquotedblleft{}minimum; at least\textquotedblright{}.

\begin{quote}
\mr{या} \mr{कामासाठी} \mr{किमान} \mr{दोन} \mr{दिवस} \mr{लागतील}.
\end{quote}

\begin{quote}
\emph{yā kāmāsāṭhī kimān don divas lāgtīl.}
\end{quote}

\begin{quote}
This work will take at least two days.
\end{quote}

Before a number, at least: \mr{किमान} \mr{दोन}, at least two.

\begin{cousinweb}
\textbf{\mr{किमान}} (\emph{kimān}) is a Marathi word whose origin is not settled.
\end{cousinweb}

\begin{grammarlens}[title={What you've built}]
One more word for saying what happens if.
\end{grammarlens}

\subsection*{Your turn}
\begin{itemize}
\raggedright
  \item \emph{Say it:} \emph{kimān}
  \item \emph{Say it:} \emph{kimān}, once more
  \item \emph{Say it:} \emph{yā kāmāsāṭhī kimān don divas lāgtīl.} — This work will take at least two days.
  \item \emph{Recall:} say the Marathi for to fulfil, then the Marathi for a limit, then say \emph{kimān} again
\end{itemize}

\begin{tcolorbox}[breakable,colback=teal!4,colframe=teal!35!black,title={Before you move on}]
What does \mr{किमान} mean? (\textquotedblleft{}Minimum; at least\textquotedblright{}.) Say it once more.
\end{tcolorbox}
\section[\texorpdfstring{\mr{कमाल} (kamāl)}{कमाल (kamāl)}]{\mr{कमाल} (kamāl) — maximum}

\label{lesson:MR-C379-kamal}



\begin{quote}
Before the new one: say the Marathi for a limit, then the Marathi for minimum.
\end{quote}

\subsection*{You'll want to know: \mr{कमाल}}
\textbf{\mr{कमाल}} — \emph{kamāl} — \textquotedblleft{}maximum\textquotedblright{}.

\begin{quote}
\mr{आज} \mr{कमाल} \mr{तापमान} \mr{चाळीस} \mr{अंश} \mr{होतं}.
\end{quote}

\begin{quote}
\emph{āj kamāl tāpmān cāḷīs aṁś hotaṁ.}
\end{quote}

\begin{quote}
Today's maximum temperature was forty degrees.
\end{quote}

In speech also amazing: \mr{कमाल} \mr{आहे}!

\begin{cousinweb}
\textbf{\mr{कमाल}} (\emph{kamāl}) is a loan from Persian, usually traced to Arabic \emph{kamāl}, \textquotedblleft{}perfection, completeness\textquotedblright{}: the fullest amount.
\end{cousinweb}

\begin{grammarlens}[title={What you've built}]
One more word for saying what happens if.
\end{grammarlens}

\subsection*{Your turn}
\begin{itemize}
\raggedright
  \item \emph{Say it:} \emph{kamāl}
  \item \emph{Say it:} \emph{kamāl}, once more
  \item \emph{Say it:} \emph{āj kamāl tāpmān cāḷīs aṁś hotaṁ.} — Today's maximum temperature was forty degrees.
  \item \emph{Recall:} say the Marathi for a limit, then the Marathi for minimum, then say \emph{kamāl} again
\end{itemize}

\begin{tcolorbox}[breakable,colback=teal!4,colframe=teal!35!black,title={Before you move on}]
What does \mr{कमाल} mean? (\textquotedblleft{}Maximum\textquotedblright{}.) Say it once more.
\end{tcolorbox}
\section[\texorpdfstring{\mr{शिथिल} \mr{करणे} (śithil karṇe)}{शिथिल करणे (śithil karṇe)}]{\mr{शिथिल} \mr{करणे} (śithil karṇe) — to relax (a rule)}

\label{lesson:MR-C379-sithil-karne}



\begin{quote}
Before the new one: say the Marathi for minimum, then the Marathi for maximum.
\end{quote}

\subsection*{You'll want to know: \mr{शिथिल} \mr{करणे}}
\textbf{\mr{शिथिल} \mr{करणे}} — \emph{śithil karṇe} — \textquotedblleft{}to relax (a rule)\textquotedblright{}.

\begin{quote}
\mr{पाऊस} \mr{थांबल्यावर} \mr{प्रशासनाने} \mr{निर्बंध} \mr{शिथिल} \mr{केले}.
\end{quote}

\begin{quote}
\emph{pāūs thāmbalyāvar praśāsnāne nirbandh śithil kele.}
\end{quote}

\begin{quote}
When the rain stopped, the administration relaxed the restrictions.
\end{quote}

Past, with -\mr{ने} (he …): \mr{त्याने} \mr{नियम} \mr{शिथिल} \mr{केला}, \emph{tyāne niyam śithil kelā}.

\mr{शिथिल}, slack, loose; \mr{शिथिल} \mr{होणे}, to slacken.

\begin{cousinweb}
\textbf{\mr{शिथिल} \mr{करणे}} (\emph{śithil karṇe}) joins the tatsama \mr{शिथिल}, Sanskrit \emph{śithila}, \textquotedblleft{}loose, slack\textquotedblright{} (Monier-Williams), to \mr{करणे}, \textquotedblleft{}to make\textquotedblright{}.
\end{cousinweb}

\begin{grammarlens}[title={What you've built}]
One more word for saying what happens if.
\end{grammarlens}

\subsection*{Your turn}
\begin{itemize}
\raggedright
  \item \emph{Say it:} \emph{śithil karṇe}
  \item \emph{Say it:} \emph{śithil karṇe}, once more
  \item \emph{Say it:} \emph{pāūs thāmbalyāvar praśāsnāne nirbandh śithil kele.} — When the rain stopped, the administration relaxed the restrictions.
  \item \emph{Recall:} say the Marathi for minimum, then the Marathi for maximum, then say \emph{śithil karṇe} again
\end{itemize}

\begin{tcolorbox}[breakable,colback=teal!4,colframe=teal!35!black,title={Before you move on}]
What does \mr{शिथिल} \mr{करणे} mean? (\textquotedblleft{}To relax (a rule)\textquotedblright{}.) Say it once more.
\end{tcolorbox}
\section[\texorpdfstring{\mr{झेपणे} (jhepṇe)}{झेपणे (jhepṇe)}]{\mr{झेपणे} (jhepṇe) — to be manageable (for someone)}

\label{lesson:MR-C379-jhepne}



\begin{quote}
Before the new one: say the Marathi for maximum, then the Marathi for to relax.
\end{quote}

\subsection*{You'll want to know: \mr{झेपणे}}
\textbf{\mr{झेपणे}} — \emph{jhepṇe} — \textquotedblleft{}to be manageable (for someone)\textquotedblright{}.

\begin{quote}
\mr{एवढं} \mr{काम} \mr{मला} \mr{एकट्याला} \mr{झेपणार} \mr{नाही}.
\end{quote}

\begin{quote}
\emph{evḍhaṁ kām malā ekaṭyālā jhepṇār nāhī.}
\end{quote}

\begin{quote}
I won't be able to manage so much work alone.
\end{quote}

Past (to him …): \mr{त्याला} \mr{झेपलं}, \emph{tyālā jheplaṁ}.

Dative: \mr{मला} \mr{झेपतं}, I can cope; mostly with \mr{नाही}.

\begin{cousinweb}
\textbf{\mr{झेपणे}} (\emph{jhepṇe}) is a Marathi verb whose origin is not settled.
\end{cousinweb}

\begin{grammarlens}[title={What you've built}]
One more word for saying what happens if.
\end{grammarlens}

\subsection*{Your turn}
\begin{itemize}
\raggedright
  \item \emph{Say it:} \emph{jhepṇe}
  \item \emph{Say it:} \emph{jhepṇe}, once more
  \item \emph{Say it:} \emph{evḍhaṁ kām malā ekaṭyālā jhepṇār nāhī.} — I won't be able to manage so much work alone.
  \item \emph{Recall:} say the Marathi for maximum, then the Marathi for to relax, then say \emph{jhepṇe} again
\end{itemize}

\begin{tcolorbox}[breakable,colback=teal!4,colframe=teal!35!black,title={Before you move on}]
What does \mr{झेपणे} mean? (\textquotedblleft{}To be manageable (for someone)\textquotedblright{}.) Say it once more.
\end{tcolorbox}
